\documentclass[oneside,openany,headings=optiontotoc,11pt,numbers=noenddot]{scrreprt}

\usepackage[a4paper]{geometry}
\usepackage[utf8]{inputenc}
\usepackage[T1]{fontenc}
\usepackage{lmodern}
\usepackage[ngerman]{babel}
\usepackage{ngerman}

\usepackage[onehalfspacing]{setspace}

\usepackage{fancyhdr}
\usepackage{fancybox}

\usepackage{rotating}
\usepackage{varwidth}

%Struktogramme
\usepackage[german,curves]{struktex}

\usepackage{pdflscape}
\usepackage{changepage}
\usepackage{graphicx}
\usepackage[bottom]{footmisc}
\usepackage{transparent}
\usepackage{graphbox}
\graphicspath{
	{Pics/PDFs/}
	{Pics/JPGs/}
	{Pics/PNGs/}
}
\usepackage{caption}
\usepackage{wrapfig}
\usepackage{marginnote}
\usepackage{tabularx}
\usepackage{dashrule}
\usepackage{soulutf8}
\usepackage{hhline}
%arydshln suppresses vertical lines in table
%\usepackage{arydshln}
\usepackage{multirow}
\usepackage{enumerate}
\usepackage[hidelinks]{hyperref}
\usepackage{listings}

\usepackage[table]{xcolor}
\usepackage{array}
\usepackage{enumitem,amssymb,amsmath}
\usepackage{interval}
\usepackage{cancel}
\usepackage{stmaryrd}
\usepackage{wasysym}
\usepackage{polynom}
\usepackage{diagbox}
\usepackage{dashrule}
\usepackage{framed}
\usepackage{mdframed}
\usepackage{karnaugh-map}
\usepackage{pdfpages}

\usepackage{blindtext}

\usepackage{eso-pic}

\usepackage{amssymb}
\usepackage{eurosym}

\usepackage[pages=some]{background}
\pagestyle{headings}
\renewcommand{\headrulewidth}{0.2pt}
\renewcommand{\footrulewidth}{0.2pt}
\newcommand*{\underdownarrow}[2]{\ensuremath{\underset{\overset{\Big\downarrow}{#2}}{#1}}}
\setlength{\fboxsep}{5pt}
\newcommand{\explainBelow}[3]{\underbrace{#1}_{\parbox{\widthof{#3}}{\footnotesize\raggedright #2}}}
\newcommand{\explainAbove}[3]{\overbrace{#1}^{\parbox{\widthof{#3}}{\footnotesize\raggedright #2}}}
\newcommand\footnoteref[1]{\protected@xdef\@thefnmark{\ref{#1}}\@footnotemark}


% Codestyle defined
\definecolor{codegreen}{rgb}{0,0.6,0}
\definecolor{codegray}{rgb}{0.5,0.5,0.5}
\definecolor{codepurple}{rgb}{0.58,0,0.82}
\definecolor{backcolour}{rgb}{0.95,0.95,0.92}
\definecolor{deepgreen}{rgb}{0,0.5,0}
\definecolor{darkblue}{rgb}{0,0,0.65}
\definecolor{mauve}{rgb}{0.40, 0.19,0.28}
\colorlet{exceptioncolour}{yellow!50!red}
\colorlet{commandcolour}{blue!60!black}
\colorlet{numpycolour}{blue!60!green}
\colorlet{specmethodcolour}{violet}

%Neue Spaltendefinition
\newcolumntype{L}[1]{>{\raggedright\let\newline\\\arraybackslash\hspace{0pt}}m{#1}}
\newcolumntype{M}{>{\centering\arraybackslash}X}
\newcommand{\cmnt}[1]{\ignorespaces}
%Textausrichtung ändern
\newcommand\tabrotate[1]{\rotatebox{90}{\raggedright#1\hspace{\tabcolsep}}}

%Intervall-Konfig
\intervalconfig {
	soft open fences
}

%Bash
\lstdefinestyle{BashInputStyle}{
	language=bash,
	basicstyle=\small\sffamily,
	backgroundcolor=\color{backcolour},
	columns=fullflexible,
	backgroundcolor=\color{backcolour},
	breaklines=true,
}
%Java
\lstdefinestyle{JavaInputStyle}{
	language=Java,
	backgroundcolor=\color{backcolour},
	aboveskip=1mm,
	belowskip=1mm,
	showstringspaces=false,
	columns=flexible,
	basicstyle={\footnotesize\ttfamily},
	numberstyle={\tiny},
	numbers=none,
	keywordstyle=\color{purple},,
	commentstyle=\color{deepgreen},
	stringstyle=\color{blue},
	emph={out},
	emphstyle=\color{darkblue},
	emph={[2]rand},
	emphstyle=[2]\color{specmethodcolour},
	breaklines=true,
	breakatwhitespace=true,
	tabsize=2,
}
%Python
\lstdefinestyle{PythonInputStyle}{
	language=Python,
	alsoletter={1234567890},
	aboveskip=1ex,
	basicstyle=\footnotesize,
	breaklines=true,
	breakatwhitespace= true,
	backgroundcolor=\color{backcolour},
	commentstyle=\color{red},
	otherkeywords={\ , \}, \{, \&,\|},
	emph={and,break,class,continue,def,yield,del,elif,else,%
		except,exec,finally,for,from,global,if,import,in,%
		lambda,not,or,pass,print,raise,return,try,while,assert},
	emphstyle=\color{exceptioncolour},
	emph={[2]True,False,None,min},
	emphstyle=[2]\color{specmethodcolour},
	emph={[3]object,type,isinstance,copy,deepcopy,zip,enumerate,reversed,list,len,dict,tuple,xrange,append,execfile,real,imag,reduce,str,repr},
	emphstyle=[3]\color{commandcolour},
	emph={[4]ode, fsolve, sqrt, exp, sin, cos, arccos, pi,  array, norm, solve, dot, arange, , isscalar, max, sum, flatten, shape, reshape, find, any, all, abs, plot, linspace, legend, quad, polyval,polyfit, hstack, concatenate,vstack,column_stack,empty,zeros,ones,rand,vander,grid,pcolor,eig,eigs,eigvals,svd,qr,tan,det,logspace,roll,mean,cumsum,cumprod,diff,vectorize,lstsq,cla,eye,xlabel,ylabel,squeeze},
	emphstyle=[4]\color{numpycolour},
	emph={[5]__init__,__add__,__mul__,__div__,__sub__,__call__,__getitem__,__setitem__,__eq__,__ne__,__nonzero__,__rmul__,__radd__,__repr__,__str__,__get__,__truediv__,__pow__,__name__,__future__,__all__},
	emphstyle=[5]\color{specmethodcolour},
	emph={[6]assert,range,yield},
	emphstyle=[6]\color{specmethodcolour}\bfseries,
	emph={[7]Exception,NameError,IndexError,SyntaxError,TypeError,ValueError,OverflowError,ZeroDivisionError,KeyboardInterrupt},
	emphstyle=[7]\color{specmethodcolour}\bfseries,
	emph={[8]taster,send,sendMail,capture,check,noMsg,go,move,switch,humTem,ventilate,buzz},
	emphstyle=[8]\color{blue},
	keywordstyle=\color{blue}\bfseries,
	rulecolor=\color{black!40},
	showstringspaces=false,
	stringstyle=\color{deepgreen}
}

\lstset{literate=%
	{Ö}{{\"O}}1
	{Ä}{{\"A}}1
	{Ü}{{\"U}}1
	{ß}{{\ss}}1
	{ü}{{\"u}}1
	{ä}{{\"a}}1
	{ö}{{\"o}}1
}

% Neue Klassenarbeits-Umgebung
\newenvironment{worksheet}[3]
% Begin-Bereich
{
	\newpage
	\sffamily
	\setcounter{page}{1}
	\ClearShipoutPicture
	\AddToShipoutPicture{
		\put(55,761){{
				\mbox{\parbox{385\unitlength}{\tiny \color{codegray}BBS I Mainz, #1 \newline #2
						\newline #3
					}
				}
			}
		}
		\put(455,761){{
				\mbox{\hspace{0.3cm}\includegraphics[width=0.2\textwidth]{../../logo.pdf}}
			}
		}
	}
}
% End-Bereich
{
	\clearpage
	\ClearShipoutPicture
}

\geometry{left=2.50cm,right=2.50cm,top=3.00cm,bottom=1.00cm,includeheadfoot}

\begin{document}
	\begin{worksheet}{Höhere Berufsfachschule IT-Systeme}{Grundstufe - Mathematik}{Wochenplan Lineare Funktionen}
		\noindent
		\begin{tabularx}{\textwidth}{XXl}
			Wochenplan Nr.: \rule{0.15\textwidth}{1pt} & Erledigt: & Zeitraum: \underline{10.09 - 16.09}
		\end{tabularx}
	
		\begin{framed}
			\noindent
			\textbf{Montag:} Geben Sie die Steigung und den y-Achsenabschnitt der linearen Funktion an.\\
			\par\noindent
			\begin{tabularx}{\textwidth}{XX}
				(a) \(f(x) = 3x\) & (b) \(f(x) = \frac{1}{5}x +2\)\\
				\colorbox{green!10}{\(m = 3; y-AAS = 0\)} & \colorbox{green!10}{\(m = \frac{1}{5}; y-AAS = 2\)}\\
				\\
				(c) \(f(x) = -\frac{4}{3}x -\frac{5}{2}\) & (d) \(f(x) = 1,5x+0,5\)\\
				\colorbox{green!10}{\(m=-\frac{4}{3}; y-AAS = -\frac{5}{2}\)} & \colorbox{green!10}{\(m = 1,5; y-AAS = 0,5\)}
			\end{tabularx}
		\end{framed}
		\begin{framed}
			\noindent
			\textbf{Dienstag:} Bestimmen Sie die Funktionsgleichung der wie folgt gegebenen linearen Funktion.
			\indent{} Die Gerade steigt um ein Drittel pro Einheit auf der x-Achse und geht durch den Punkt \(P(-3|-4)\).
			\begin{center}
				\rule{0.9\textwidth}{0.1pt}\\
			\end{center}
			Wir erinnern uns, dass eine lineare Funktion die Form \colorbox{green!10}{\(y = mx + b\)} hat, wobei \underline{m die Steigung} und \underline{b den y-Achsenabschnitt} angibt.\\
			\par\noindent
			Aus dem Text wissen wir: \(m = \frac{1}{3}\). Wir benötigen also lediglich noch den y-Achsenabschnitt. Um diesen zu bestimmen, verwenden wir die \textbf{Punkt-Steigung-Form} \(y-y_1 = m*(x-x_1)\).\\
			\par\noindent
			Der gegebene Punkt \(P(\underbrace{-3}_{x_1}|\underbrace{-4}_{y_1})\) gibt uns die nötigen Informationen. Wir setzen diese nur noch in die Gleichung ein und lösen anschließend nach \(y\) auf.\\
			\par\noindent
			\begin{tabularx}{\textwidth}{ll}
				\(y - (-4) = \frac{1}{3}\cdot(x-(-3))\) & |\(AM\)\\
				\(y + 4 = \frac{1}{3}x + 1\) & |\(-4\)\\
				\colorbox{green!10}{\(y = \frac{1}{3}x -1\)}
			\end{tabularx}
		\end{framed}
		\begin{framed}
			\noindent
			\textbf{Mittwoch:} Bestimmen Sie jeweils die Gleichung zu der Geraden, die durch P geht und die Steigung \(m\) hat.
			\begin{framed}
				\noindent
				Wir kennen sowohl einen Punkt wie auch die Steigung der Geraden. Das bedeutet, um die entsprechende Funktionsgleichung aufzustellen, benötigen wir lediglich die \textbf{Punkt-Steigung-Form}: \(\mathbf{y -y_1 = m\cdot(x-x_1)}\).
			\end{framed}
			\begin{tabularx}{\textwidth}{Xl|Xl}
				(a) \(P(2|5); m=3\) & & (b) \(P(4|-2); m=0\)\\
				\(y - 5 = 3(x-2)\) & |\(AM\) & \(y-(-2) = 4(x-4)\) & |\(AM\)\\
				\(y - 5 = 3x -6\) & |\(+5\) & \(y+2 = 4x -16\) & |\(-2\)\\
				\colorbox{green!10}{\(y= 3x - 1\)} & & \colorbox{green!10}{\(y = 4x - 18\)}\\
				\\
				\hline
				\\
				(c) \(P(-3|1); m=-1\) & & (d) \(P(1,5|0,5); m=-4\)\\
				\(y - 1 = (-1)(x-(-3))\) &  & \( y-0,5 = (-4)(x-1,5)\) & |\(AM\)\\
				\(y -1 = (-1)(x+3)\) & |\(AM\) & \( y-0,5 = -4x +6\) & |\(+0,5\)\\
				\(y - 1 = -x -3\) & |\(+1\) & \colorbox{green!10}{\(y = -4x +6,5\)}\\
				\colorbox{green!10}{\(y = -x -2\)} & & 
			\end{tabularx}
		\end{framed}
		\begin{framed}
			\noindent
			\textbf{Donnerstag:} Käpt\grq{}n Blaubär fährt mit seinem Tankschiff A bei einer Durchschnittsgeschwindigkeit von \(400\ sm\) pro Tag von Hong Kong nach Hamburg.\\
			Hein Blöd steuert Tankschiff B von Hamburg nach Hong Kong mit \(550\ sm\) pro Tag zurück.
			\begin{itemize}
				\item[(b)] Geben Sie die Funktionsgleichung von \(f_A\) bzw. \(f_B\) an, die die Fahrt der Tanker A bzw B beschreiben.\\
				\colorbox{green!10}{\(f_A(x) = 400x; f_B(x) = 550x\)}
				\item[(c)] Stellen Sie zu jeder Funktion eine Wertetabelle auf.
			\end{itemize}
			\begin{tabularx}{\textwidth}{l|l|l|l|l|l|l|l|l}
				x & 0 & 1 & 2 & 3 & 4 & 5 & 6 & 7\\
				\hline
				& & & & & & & &\\
				\(f_A(x)\) & 0 & 400 & 800 & 1200 & 1600 & 2000 & 2400 & 2800\\
				\hline
				& & & & & & & &\\
				\(f_B(x)\) & 0 & 550 & 1100 & 1650 & 2200 & 2750 & 3300 & 3850\\
			\end{tabularx}
			\begin{itemize}
				\item[(d)] Nach wie vielen Tagen können  sich Käpt\grq{}n Blaubär und Hein Blöd auf hoher See zuwinken?\\
				\colorbox{red!10}{Da wir keine Information über die Strecke haben, die zwischen Hong Kong und Ham-}\\
				\colorbox{red!10}{burg liegt, wir also nicht sagen können, wie weit die beiden bereits von ihrem Ursprungs-}\\
				\colorbox{red!10}{ort entfernt sind, lässt sich die Frage mit den gegebenen Informationen nicht lösen.}
			\end{itemize}
		\end{framed}
		\begin{framed}
			\noindent
			\textbf{Freitag:} Zwei Motorradfahrer fahren auf derselben Straße von A nach B. Die beiden Orte sind \(270\ km\) voneinander entfernt.\\
			Fahrer M1 fährt um 9 Uhr ab und hält eine Durchschnittsgeschwindigkeit von \(45\ \frac{km}{h}\). 75 Minuten später startet Fahrer M2 und fährt durchschnittlich \(60\ \frac{km}{h}\).
			\begin{itemize}
				\item[(a)] Stellen Sie den Sachverhalt mithilfe zweier Funktionen dar.\\
				Die aufgestellten Funktionen beschreiben jeweils die noch zu fahrende Strecke der Motorradfahrer in Abhängigkeit der verstrichenen Zeit seit Reisebeginn.\\
				\colorbox{green!10}{\(f_{M1}(x) = 270 - 45x\)}\\
				\colorbox{green!10}{\(f_{M2}(x) = 270 - 60x\)}\\
				\item[(b)] Ermitteln Sie durch Rechnung die Ankunftszeiten der beiden Fahrer.\\
				Hierfür müssen wir bestimmen, wieviel Zeit verstreicht, bis sie die Stadt B erreichen. Also genauer, wann die noch zu fahrende Strecke \(0\) ist.\\
				\par\noindent
				\begin{tabularx}{0.8\textwidth}{Xl|Xl}
					\(270 - 45x = 0\) & |\(+45x\) & \(270 - 60x = 0\) & |\(+60x\)\\
					\(270 = 45x\) & |\(:45\) & \(270 = 60x\) & |\(:60\)\\
					\(x = 6\) & & \(x=4,5\)\\
					\\
					\hline
					\\
				\end{tabularx}
				\begin{tabularx}{0.8\textwidth}{L{5.7cm}|L{5.7cm}}
					M1 erreicht sein Ziel nach 6 Stunden. & M2 erreicht sein Ziel bereits nach 4,5 Stunden also 4 Stunden und 30 Minuten.\\
					$\Rightarrow$ 9:00 Uhr + 6:00 Stunden = \colorbox{green!10}{15:00 Uhr} & $\Rightarrow$ 10:15 Uhr + 4:30 Stunden = \colorbox{green!10}{14:45 Uhr}
				\end{tabularx}
				\item[(c)] Zu welchem Zeitpunkt treffen sich die beiden Fahrer? Wie weit sind sie zu diesem Zeitpunkt vom Startpunkt entfernt?\\
				Gefragt ist nach dem Zeitpunkt, an dem beide Motorradfahrer die gleiche Strecke zurückgelegt haben.\\
				Hier ist zu beachten, dass M1 zum Startzeitpunkt von M2 bereits 75 Minuten unterwegs ist. Bevor man also \(f_{M1}(x) = f_{M2}(x)\) setzen kann, muss die Funktionsgleichung von M1 angepasst werden.\\
				\(\Rightarrow f_{M1}(x) = 270 - 45x - \underbrace{45*1,25}_{Gefahrene\ Strecke\ nach\ 75\ Minuten} = 270 - 45x - 56,25\)\\
				\(\Rightarrow f_{M1}(x) = 213,75 - 45x\)\\
				\\
				Nun setzen wir die Funktionen gleich und lösen nach \(x\).\\
				\begin{tabularx}{0.8\textwidth}{Xl}
					\(213,75 - 45x = 270 - 60x\) & |\(+60x; -213,75\)\\
					\(15x = 56,25\) & |\(:15\)\\
					\(x = 3,75\)
				\end{tabularx}\\
				\\
				Da nach dem Zeitpunkt gefragt ist, addieren wir wie oben\\
				10:15 Uhr + 3:45 Stunden = \colorbox{green!10}{14:00 Uhr}\\
				\\
				Ebenso fehlt uns noch die Information, wie weit sie zu diesem Zeitpunkt vom Startpunkt entfernt sind. Dafür müssen wir lediglich bestimmen, wie weit sie in 3,75 Stunden gefahren sind.\\
				\(\Rightarrow 60*3,75 =\) \colorbox{green!10}{\(225 km\)}.
			\end{itemize}
		\end{framed}
	\end{worksheet}
\end{document}